\section{related work}
\noindent \textbf{Hard Prompt Compression for LLMs.}
Hard prompt compression improves LLM efficiency by removing redundant content from prompts \citep{DBLP:conf/naacl/LiLSC25}. 
Methods recently have evolved from using smaller language models \citep{DBLP:conf/emnlp/JiangWLYQ23, DBLP:conf/emnlp/0001DGL23, DBLP:conf/naacl/ChuangXCLCH24} to query-aware approaches that preserve task-relevant information \citep{weston2023system, DBLP:conf/iclr/CreswellSH23, DBLP:conf/acl/JiangWL0L0Q24}. 
Additionally, lightweight bidirectional encoders have demonstrated strong effectiveness and efficiency \citep{DBLP:conf/acl/PanWJXLZLR0LZQ024, DBLP:conf/aaai/LiskavetsURKEL25}.

\noindent \textbf{Chunking Strategies in RAG Systems.}
Retrieval-Augmented Generation (RAG) systems rely on chunking extrernal documents into smaller units for retrieval \citep{DBLP:conf/nips/LewisPPPKGKLYR020, gao2023retrieval}. 
Existing chunking strategies include rule-based methods creating fixed-size segments \citep{DBLP:conf/nips/LewisPPPKGKLYR020, DBLP:conf/iclr/SarthiATKGM24, Edge2024FromLT, DBLP:conf/nips/GutierrezS0Y024}, semantic-based methods grouping content by topic \citep{DBLP:conf/naacl/QuTB25}, and LLM-driven methods leveraging model knowledge for splitting \citep{DBLP:conf/iclr/PanWJLCL0LZQ025, DBLP:conf/emnlp/DuarteMGF0O24, Zhao2024MetaChunkingLT, DBLP:conf/ecir/LiuSC25}. 
However, all of these chunking strategies for RAG systems are tailored to static scenarios, not applicable to dynamic and open-ended environments.
% Version 2 
% RAG systems enhance LLMs by retrieving relevant document chunks from external sources based on queries \citep{DBLP:conf/nips/LewisPPPKGKLYR020}. The first step in the RAG pipeline is indexing, which splits documents into smaller units \citep{gao2023retrieval}. Existing chunking strategies can be broadly categorized into rule-based, semantic-based, and LLM-driven approaches. Rule-based chunking divides documents into fixed-size segments, either with or without overlaps. This strategy is widely adopted in many RAG frameworks \citep{DBLP:conf/nips/LewisPPPKGKLYR020, DBLP:conf/iclr/SarthiATKGM24, Edge2024FromLT, DBLP:conf/nips/GutierrezS0Y024}. LangChain \citep{langchain} implements a more advanced version, recursively splitting text based on a predefined list of separators. Semantic-based chunking groups documents into semantically coherent segments, including breakpoint-based or clustering-based semantic chunking \citep{DBLP:conf/naacl/QuTB25}. LLM-driven methods leverage general knowledge from LLMs to complete the splitting process. For instance, SeCom \citep{DBLP:conf/iclr/PanWJLCL0LZQ025} and LumberChunker \citep{DBLP:conf/emnlp/DuarteMGF0O24} prompt LLMs to segment documents. By contrast, \citet{Zhao2024MetaChunkingLT} and \citet{DBLP:conf/ecir/LiuSC25} split texts based on logits computed by LLMs. However, all of these chunking strategies for RAG systems are tailored to static scenarios, not applicable to dynamic and open-ended environments.
% Version 1 
% RAG systems enhance LLMs by retrieving relevant document chunks from external sources based on queries \citep{DBLP:conf/nips/LewisPPPKGKLYR020}. The first step in the RAG pipeline is indexing, which splits documents into smaller units \citep{gao2023retrieval}. Existing chunking strategies can be broadly categorized into rule-based, semantic-based, and LLM-driven approaches. Rule-based chunking divides documents into fixed-size segments, either with or without overlaps. This strategy is widely adopted in many RAG frameworks \citep{DBLP:conf/nips/LewisPPPKGKLYR020, DBLP:conf/iclr/SarthiATKGM24, Edge2024FromLT, DBLP:conf/nips/GutierrezS0Y024}. LangChain \citep{langchain} implements a more advanced version, recursively splitting text based on a predefined list of separators. By contrast, semantic-based chunking groups documents into semantically coherent segments, including breakpoint-based or clustering-based semantic chunking \citep{DBLP:conf/naacl/QuTB25}. LLM-driven methods leverage general knowledge from LLMs to guide the splitting process. For instance, SeCom \citep{DBLP:conf/iclr/PanWJLCL0LZQ025} uses GPT-4 for zero-shot chunking. Meta-Chunking \citep{Zhao2024MetaChunkingLT} calculates the perplexity of each sentence based on its preceding context and segments documents accordingly. Different from Meta-Chunking, \citet{DBLP:conf/ecir/LiuSC25} estimate the likelihood of an [EOS] (End of Sequence) token after each sentence. However, all of these methods are designed for indexing scenarios in RAG, not applicable to dynamic and open-ended environments.

\noindent \textbf{Memory Systems for LLM Agents.}
% Too Long
% Memory systems help LLM agents move beyond stateless interactions to support flexible reasoning and adaptation in complex and changing environments \citep{liu2025advances, mei2025survey, Shan2025CognitiveMI, wu2025human}.
Memory systems help LLM agents move beyond stateless interactions to support flexible reasoning and adaptation in complex and changing environments \citep{liu2025advances, mei2025survey}. The earliest and most straightforward approaches store experiences as linear or sequential streams, sometimes enhanced with hierarchical structures \citep{Liang2023SCMEL, DBLP:conf/uist/ParkOCMLB23, Packer2023MemGPTTL, DBLP:conf/aaai/ZhongGGYW24, Salama2025MemInsightAM,fang2025mempexploringagentprocedural}. A more structured class of methods represents memories as nodes and their relationships as edges, using trees, graphs, or temporal knowledge structures to support retrieval and update \citep{DBLP:conf/iclr/RezazadehLWB25, Chhikara2025Mem0BP, rasmussen2025zep, Xu2025AMEMAM, zhang2025g}. The latest trend integrates various types of memory, allowing them to interact and synergistically improve overall performance \citep{Kang2025MemoryOO, Li2025MemOSAO, wang2025mirix, nan2025nemori}. Overall, existing memory systems for LLM agents have become increasingly complex and capable, leveraging hierarchical, structured, and multi-type memories. However, most focus on maximizing effectiveness, with limited consideration of efficiency. While some recent works \citep{guo2024lightrag, zhao20252graphrag, dong2025youtu} share a similar motivation with our work, they focus on lightweight adaptations of GraphRAG where the corpus is predefined and static. 

% Version 1
% Memory systems help LLM agents move beyond stateless interactions by allowing them to build on past experiences, which, in turn, supports more flexible reasoning and adaptation in complex and changing environments \citep{liu2025advances, mei2025survey, Shan2025CognitiveMI, wu2025human}. \citet{Liang2023SCMEL} proposed the SCM framework for LLM agents, which includes a memory stream stored in an external database for conversation histories and a memory controller that determines how to use the memory stream. \citet{DBLP:conf/uist/ParkOCMLB23} considered recency, importance, and relevance of memories during retrieval. MemoryBank \citep{DBLP:conf/aaai/ZhongGGYW24} focused on personalization, introducing a forgetting mechanism based on the Ebbinghaus forgetting curve. \citet{Packer2023MemGPTTL} drew inspiration from hierarchical memory systems in operating systems, introducing a two-level memory structure (main context and external context) with paging between them. \citet{DBLP:conf/iclr/RezazadehLWB25} organized memory as a tree structure and proposed a memory update algorithm tailored to tree-structured representations. \citet{Chhikara2025Mem0BP} presented the Mem0 framework, which includes conflict detection and organizes memory as a graph. \citet{rasmussen2025zep} organized agent memory with a temporal knowledge graph. Inspired by the structuring principles of Zettelkasten, \citet{Xu2025AMEMAM} built the A-MEM framework, which generates links between memory units and allows prior memory units to evolve. \citet{zhang2025g} designed G-Memory with a three-tier graph hierarchy to support multi-agent systems. Recent works including MemoryOS \citep{Kang2025MemoryOO} and MemOS \citep{Li2025MemOSAO} improve memory systems by combining multiple types of memories to enhance overall performance. Overall, existing memory systems for LLM agents have become increasingly complex and capable, leveraging hierarchical, graph-based, and multi-type memories. However, most focus on maximizing effectiveness, with limited consideration of efficiency.
% In this work, we propose a lightweight memory architecture that achieves a balance between performance and efficiency, addressing this overlooked aspect of memory design. 

% MemoryBank \citep{DBLP:conf/aaai/ZhongGGYW24} focused on personalization by storing conversation logs, user characteristics, and related events, introducing a forgetting mechanism based on the Ebbinghaus forgetting curve.

% From the perspective of memory representation, memory systems can be divided into Parametric Memory Systems and Non-parametric Memory Systems. Parametric memory systems refer to approaches where information is stored directly within model parameters. A seminal contribution in this direction is the Long Short-Term Memory (LSTM) architecture \citep{DBLP:journals/neco/HochreiterS97}, which represents memory using a hidden state vector. Through training, neural networks learn to update this hidden state and retrieve information from it. However, the representational capacity of such hidden vectors is inherently limited. To address this, subsequent studies have expanded memory size \citep{DBLP:journals/corr/WestonCB14, DBLP:journals/nature/GravesWRHDGCGRA16, DBLP:conf/nips/BulatovKB22, DBLP:conf/icml/WangGCJLYYLLYSM24, berges2024memory} or introduced more sophisticated update mechanisms \citep{sun2024learning, behrouz2024titans, Wang2025MEM, DBLP:conf/nips/YangWZSK24} to improve parametric memory. Despite these advances, parameteric memory faces two persistent challenges \citep{du2025rethinking}. First, it lacks interpretability, as knowledge is entangled within model parameters. Second, it is difficult to modify, even with progress in knowledge editing \citep{zhang2024comprehensive, DBLP:journals/csur/WangZLZCL25} and machine unlearning techniques \citep{DBLP:journals/air/BlancoJusticiaJMSDCT25, liu2024machine}.
